% =====================================================================
%  Bitácora de desarrollo — Laboratorio Virtual de Estabilidad
%  Compilar con LuaLaTeX (o XeLaTeX):  lualatex bitacora.tex  (dos veces)
%  En Overleaf: Menú → Compiler → LuaLaTeX
% =====================================================================
\documentclass[11pt,a4paper]{article}

\usepackage{fontspec}
\setmainfont{DejaVuSerif}[Extension=.ttf, BoldFont=*-Bold, ItalicFont=*-Italic, BoldItalicFont=*-BoldItalic, Scale=0.92]
\setsansfont{DejaVuSans}[Extension=.ttf, BoldFont=*-Bold, ItalicFont=*-Oblique, BoldItalicFont=*-BoldOblique]
\setmonofont{DejaVuSansMono}[Extension=.ttf, BoldFont=*-Bold, ItalicFont=*-Oblique, BoldItalicFont=*-BoldOblique, Scale=0.86]
\usepackage[spanish,es-noshorthands,es-tabla]{babel}
\usepackage[margin=2.3cm,headheight=14pt]{geometry}
\usepackage{amsmath}
\usepackage{xcolor}
\usepackage{booktabs}
\usepackage{tabularx}
\usepackage{longtable}
\usepackage{array}
\usepackage{enumitem}
\usepackage{fvextra}
\usepackage{fancyhdr}
\usepackage{titlesec}
\usepackage[hidelinks]{hyperref}

\definecolor{acento}{HTML}{0284C7}
\definecolor{estable}{HTML}{059669}
\definecolor{vuelco}{HTML}{DC2626}
\definecolor{aviso}{HTML}{D97706}
\definecolor{gris}{HTML}{475569}
\definecolor{fondo}{HTML}{F1F5F9}

\titleformat{\section}{\Large\bfseries\sffamily\color{acento}}{\thesection}{0.8em}{}
\titleformat{\subsection}{\large\bfseries\sffamily}{\thesubsection}{0.8em}{}
\titleformat{\subsubsection}{\normalsize\bfseries\sffamily}{\thesubsubsection}{0.8em}{}

\pagestyle{fancy}
\fancyhf{}
\fancyhead[L]{\small\sffamily\color{gris} Bitácora · Laboratorio Virtual de Estabilidad}
\fancyhead[R]{\small\sffamily\color{gris} Sara · Hugo · Daniela}
\fancyfoot[C]{\small\sffamily\thepage}
\renewcommand{\headrulewidth}{0.4pt}

\newcolumntype{L}[1]{>{\raggedright\arraybackslash}p{#1}}
\newcolumntype{Y}{>{\raggedright\arraybackslash}X}
\setlist{itemsep=2pt,topsep=4pt}

% Bloque para citar literalmente las instrucciones (prompts)
\DefineVerbatimEnvironment{Prompt}{Verbatim}{breaklines, breakanywhere, fontsize=\footnotesize, frame=single, framesep=3mm, rulecolor=\color{acento}, numbers=none}

\newcommand{\miembro}[1]{\textbf{\sffamily #1}}
\newcommand{\ok}{\textcolor{estable}{\textbf{\sffamily ✔}}}

\begin{document}

% ---------------------------------------------------------------------
\begin{titlepage}
  \centering
  \vspace*{2.5cm}
  {\sffamily\color{acento}\Large Cálculo Integral · Aplicaciones de la integral\par}
  \vspace{1.2cm}
  {\sffamily\bfseries\Huge Bitácora de desarrollo\par}
  \vspace{0.4cm}
  {\sffamily\LARGE Laboratorio Virtual de Estabilidad\par}
  {\sffamily\large\color{gris} CAD \& Mass Properties — centroides y centros de masa\par}
  \vspace{2.2cm}
  \begin{tabular}{l}
    \toprule
    \textbf{Integrantes}\\
    \midrule
    Sara\\
    Hugo\\
    Daniela\\
    \bottomrule
  \end{tabular}
  \vfill
  {\large 20 de noviembre de 2026\par}
\end{titlepage}

\tableofcontents
\newpage

% ---------------------------------------------------------------------
\section{Objetivo del proyecto}

Construir una aplicación web de nivel profesional que sirva como laboratorio virtual para el tema
\emph{aplicaciones de la integral: centroides y centros de masa}. La aplicación permite:

\begin{itemize}
  \item modelar cuerpos planos formados por figuras (rectángulos, triángulos, círculos, semicírculos,
        cuartos de círculo, trapecios, polígonos libres y regiones definidas por ecuaciones);
  \item calcular el centroide geométrico $(\bar x_G,\bar y_G)$ y el centro de masa $(\bar x_{CM},\bar y_{CM})$
        mostrando el desarrollo paso a paso (integral $\rightarrow$ sumatoria $\rightarrow$ sustitución numérica);
  \item evaluar la estabilidad del cuerpo sobre dos apoyos A y B, calcular las reacciones y
        comprobar $\sum F_y = 0$ y $\sum\tau = 0$;
  \item simular el vuelco como cuerpo rígido cuando el centro de masa cae fuera de la base;
  \item exportar el lienzo en PNG y un reporte técnico imprimible.
\end{itemize}

% ---------------------------------------------------------------------
\section{Metodología de trabajo}

El desarrollo se realizó con un \textbf{agente de programación con inteligencia artificial} (Devin)
dirigido por el equipo mediante instrucciones escritas (\emph{prompts}). El código fuente fue generado
por el agente a partir de esas instrucciones; el trabajo del equipo consistió en:

\begin{enumerate}
  \item \textbf{Especificar}: redactar el prompt inicial con los requisitos técnicos, las convenciones
        físicas, los datos de los presets con sus resultados esperados y los criterios de aceptación.
  \item \textbf{Dirigir por fases}: autorizar cada fase («continúa») sólo después de revisarla.
  \item \textbf{Verificar}: comprobar a mano los resultados de los presets, revisar la interfaz y
        probar los casos borde.
  \item \textbf{Corregir}: detectar defectos o carencias y pedir correcciones concretas.
  \item \textbf{Documentar}: registrar en esta bitácora las instrucciones y las evidencias.
\end{enumerate}

Las instrucciones se citan literalmente en la sección~\ref{sec:evidencias} como evidencia del proceso.

% ---------------------------------------------------------------------
\section{Distribución del trabajo}

El trabajo se repartió en tres áreas de responsabilidad equivalentes. Cada integrante fue responsable
de especificar, revisar y validar los entregables de su área, y los tres participaron en la redacción del
prompt inicial y en la revisión cruzada de cada fase.

\begin{table}[h]
\centering
\small
\begin{tabularx}{\textwidth}{L{2.2cm}YYY}
\toprule
 & \miembro{Sara} & \miembro{Hugo} & \miembro{Daniela}\\
\midrule
Área & Infraestructura, interfaz general y documentación & Motor matemático y física & Lienzo CAD e interacción\\
\midrule
Fase 1 & Docker, Vite, Tailwind, GitHub Actions; tema visual & Convenciones de unidades y rangos válidos; tests del modelo & Esqueleto de paneles e inputs con sliders\\
Fase 2 & Selector de presets en la barra superior & \texttt{physics.ts}, tests Vitest, datos de presets & Panel de análisis con KaTeX (desarrollo paso a paso)\\
Fase 3 & Revisión de rendimiento y accesibilidad & Cámara, cuadrícula adaptativa y sombreado por densidad (funciones puras + tests) & Lienzo: bloques, apoyos, arrastre, marcadores y plomada\\
Fase 4 & Exportación PNG, reporte técnico y README & Simulación del vuelco (dinámica de cuerpo rígido) & Tarjetas de presets, flechas de fuerzas en el lienzo\\
Rediseño & Rediseño \emph{enterprise}: barra de título, cinta de herramientas, barra de estado & Motor de figuras geométricas y analizador de ecuaciones & Paleta con arrastrar y soltar, diálogos de ecuación y polígono, inspector\\
Correcciones & Bitácora de desarrollo & Diagnóstico del error de simulación & Verificación de la corrección en el navegador\\
\midrule
Tareas & 6 & 6 & 6\\
Participación & 33,3\,\% & 33,3\,\% & 33,3\,\%\\
\bottomrule
\end{tabularx}
\caption{Reparto uniforme de responsabilidades por fase.}
\end{table}

% ---------------------------------------------------------------------
\section{Cronología de sesiones}

\begin{longtable}{L{1.2cm}L{3.3cm}L{2.2cm}L{8.0cm}}
\toprule
\textbf{Sesión} & \textbf{Instrucción} & \textbf{Responsable} & \textbf{Resultado y evidencia}\\
\midrule
\endhead
1 & Prompt inicial (§\ref{sec:prompt-inicial}) $\rightarrow$ Fase 1 & Sara (los tres redactan) &
Infraestructura Docker + Vite + Tailwind + GitHub Actions y esqueleto de la interfaz. \texttt{npm run test}: 8/8. Build con rutas relativas y fuentes empaquetadas; HMR verificado en Docker.\\
2 & Corrección 1: «aún está lejos de ser un laboratorio…» $\rightarrow$ Fase 2 & Hugo &
Motor matemático puro, 4 presets como datos y panel KaTeX. Tests: 30/30. Los 4 presets reproducen los resultados esperados.\\
3 & Corrección 2: «continúa» $\rightarrow$ Fase 3 & Daniela &
Lienzo CAD con zoom centrado en el cursor, pan, snap, arrastre de bloques y apoyos, marcadores CG/CM y plomada. Tests: 38/38. Prueba en navegador sin errores.\\
4 & Corrección 3: «continúa» $\rightarrow$ Fase 4 & Sara y Hugo &
Simulación de gravedad, tarjetas de presets, exportación PNG, reporte imprimible y README. Tests: 46/46. Preset 4 vuelca en torno a B hasta $\theta=-26{,}57^\circ$.\\
5 & Corrección 4: «continúa hasta terminar» & Daniela &
Entrega de los archivos restantes de la Fase 4 y cierre de los 8 criterios de aceptación.\\
6 & Corrección 5: estética \emph{enterprise}, figuras por ecuaciones y arrastrar y soltar & Los tres &
Generalización del modelo a 8 tipos de figura, analizador de expresiones, regiones $g(x)\le y\le f(x)$ integradas con Simpson, rediseño completo de la interfaz. Tests: 88/88 y luego 103/103.\\
7 & Corrección 6: «le doy a simular pero no hace nada» + bitácora & Hugo, Daniela y Sara &
Diagnóstico de la causa raíz y corrección (ver §\ref{sec:sim}); redacción de esta bitácora.\\
\bottomrule
\end{longtable}

% ---------------------------------------------------------------------
\section{Evidencias: instrucciones dadas al agente}\label{sec:evidencias}

\subsection{Prompt inicial}\label{sec:prompt-inicial}

Redactado en conjunto por Sara, Hugo y Daniela. Se reproduce literalmente.

\begin{Prompt}
Desarrolla lo propuesto en el promt que te envie# ROL
Actúa como Arquitecto Frontend Principal e Ingeniero de Simulación Estructural. Construye desde cero **"Laboratorio Virtual de Estabilidad (CAD & Mass Properties)"**, una app web de nivel profesional para un curso universitario de Cálculo Integral (aplicaciones de la integral: centroides y centros de masa). Estética y precisión inspiradas en SkyCiv Structural 3D, Autodesk Fusion y Algodoo. Idioma de toda la interfaz y de las explicaciones: español.

# REGLAS DE ENTREGA (obligatorias)
- Entrega en 4 FASES. Al terminar cada fase, detente y espera mi "continúa". No adelantes fases.
- Dentro de cada fase: archivos completos, sin "...", sin "resto igual", sin placeholders ni TODOs. Si una fase no cabe en una respuesta, divídela y avísame antes de truncar nada.
- Cada archivo va en su propio bloque de código, con su ruta como título.
- Al final de cada fase incluye: (a) lista de archivos creados/modificados, (b) comandos exactos para verificarla, (c) qué criterios de aceptación de la sección 11 ya se cumplen.

Fases:
1. Infraestructura (Docker, Vite, Tailwind, GitHub Actions) + esqueleto de UI con el tema aplicado.
2. Motor matemático puro (`src/lib/physics.ts`) + tests Vitest + panel de análisis con KaTeX + presets como datos.
3. Lienzo CAD (cuadrícula, ejes, zoom/pan, bloques, apoyos, arrastre, marcadores, plomada).
4. Simulación de gravedad, tarjetas de presets, exportación y README.

# 1. CONVENCIONES (no improvises, úsalas tal cual)
- Unidades: longitud en cm, masa en kg, g = 9.81 m/s². Convierte a metros solo dentro del cálculo dinámico.
- Cada bloque: `{ id, name, b, h, m, x, y }` donde **(x, y) es la esquina inferior izquierda**. Su centro es `(x + b/2, y + h/2)`.
- Eje Y positivo hacia arriba. Los apoyos A y B están en y = 0, con `xA` y `xB` editables.
- Los bloques pueden apilarse y tocarse. Si dos bloques se **solapan** (área de intersección > 0), muestra una advertencia no bloqueante en el inspector y resáltalos con borde discontinuo ámbar. Los bloques que flotan sin contacto se permiten, pero genera un aviso.
- Rangos válidos: b, h ∈ [1, 200] cm; m ∈ [0.1, 500] kg; x, y ∈ [-300, 300] cm. Valida y limita (clamp) en inputs y al arrastrar.
- Densidad: ρᵢ = mᵢ/(bᵢ·hᵢ) en kg/cm². Para el color usa la normalización (ρ − ρmin)/(ρmax − ρmin) → opacidad de relleno de 0.25 a 0.95 sobre el tono cian del bloque, más un patrón de rayado diagonal más denso a mayor ρ. Si todas las ρ son iguales, usa opacidad 0.6 y sin rayado.
- Casos borde del motor: M = 0 o A = 0 → devuelve estado `"sin_masa"` (no NaN) y la UI muestra un aviso; `xA = xB` → base de ancho 0 → INESTABLE; si `xA > xB`, ordena internamente (xmin, xmax) pero conserva las etiquetas A/B visuales.

# 2. INFRAESTRUCTURA DOCKER Y GITHUB PAGES
- `Dockerfile` con `node:20-alpine`, `WORKDIR /app`, copia `package*.json`, `npm install`, copia el resto, `EXPOSE 5173`, `CMD ["npm","run","dev","--","--host","0.0.0.0","--port","5173"]`.
- `docker-compose.yml`: servicio `app`, `build: .`, puerto `5173:5173`, volumen bind `.:/app` y volumen anónimo `/app/node_modules`, variable `CHOKIDAR_USEPOLLING=true` (para que el HMR funcione en Windows/macOS), y `vite.config.ts` con `server.watch.usePolling` coherente.
- Debe funcionar sin instalar nada en el anfitrión:
  - `docker compose up` → previsualización en http://localhost:5173
  - `docker compose run --rm app npm run build` → genera `dist/` con `index.html` en la raíz.
  - `docker compose run --rm app npm run test` → corre Vitest.
- `vite.config.ts`: `base: './'`. Todos los assets (incluidas las fuentes de KaTeX) deben salir con rutas relativas. Prohíbe CDNs externas: Inter y JetBrains Mono se empaquetan con `@fontsource/inter` y `@fontsource/jetbrains-mono`.
- `.github/workflows/deploy.yml`: en push a `main`, `actions/checkout`, `actions/setup-node` (20), `npm ci`, `npm run build`, `actions/upload-pages-artifact` (path `dist`) y `actions/deploy-pages`, con permisos `pages: write` e `id-token: write`. Explica en el README que primero hay que generar y commitear `package-lock.json` usando Docker, y activar Pages → "GitHub Actions" en el repositorio.

# 3. STACK
React 18 + TypeScript (strict) + Vite + Tailwind CSS + `lucide-react` + `katex` (importa `katex/dist/katex.min.css` en el bundle; renderiza con `katex.renderToString` y `throwOnError:false`) + Vitest. Sin otras librerías de estado o de canvas: estado con `useReducer` + Context, canvas con la API 2D nativa.

# 4. ESTÉTICA Y LAYOUT
- Colores: fondo `#0B0F19`, paneles `#111827`, bordes `#1E293B`, acento cian `#38BDF8`, estable `#10B981`, vuelco `#EF4444`, advertencia `#F59E0B`. Todos definidos como variables CSS/Tailwind.
- Tipografía: Inter para UI; JetBrains Mono con `font-variant-numeric: tabular-nums` para coordenadas, cotas y cálculos.
- Layout de workstation a pantalla completa: barra superior + panel izquierdo (320 px) + lienzo flexible + panel derecho (380 px). Paneles con scroll interno. En pantallas < 1024 px, paneles laterales colapsables en drawers.
- Barra superior: título, selector de presets, "Simular Gravedad", "Reiniciar", interruptor "Snap to Grid" (paso de 1 cm; con Shift se desactiva temporalmente), Zoom +/−, "Centrar plano", "Exportar".
- Panel izquierdo: lista de bloques (añadir, duplicar, eliminar, seleccionar; el activo resaltado) con inputs numéricos + sliders para b, h, m, x, y; sección de apoyos A y B (inputs y sliders de x).
- Panel derecho: tarjeta de diagnóstico (estado, margen, pivote), tarjeta de reacciones, y desarrollo paso a paso en KaTeX con las sustituciones numéricas.
- Accesibilidad: inputs con `label`, foco visible, `aria-live="polite"` en el diagnóstico de estabilidad, y el estado nunca depende solo del color (siempre texto + icono).

# 5. LIENZO CAD
- Canvas 2D nítido: ajusta el tamaño del buffer por `devicePixelRatio` y `ResizeObserver`.
- Cuadrícula dinámica con subdivisiones adaptativas (1, 5, 10, 50 cm según el zoom), ejes X e Y con reglas graduadas y etiquetas en mono. Zoom con la rueda centrado en el cursor (rango 0.2×–8×); pan con arrastre del botón central o Espacio + arrastre.
- Bloques: selección por clic, arrastre directo con snap opcional. Apoyos: arrastre solo horizontal. Cursor y hover coherentes.
- Símbolos: apoyo fijo (triángulo) en A y rodillo (triángulo + círculo) en B, con etiquetas.
- Base de sustentación: banda resaltada entre xmin y xmax sobre y = 0, con cota de distancia en cm.
- Marcadores diferenciados: CG (círculo con cruz, cian) y CM (círculo bicolor tipo "diana", ámbar), cada uno con etiqueta y coordenadas. Plomada: línea vertical punteada desde el CM hasta y = 0, con un punto en la proyección y la cota de distancia al apoyo más cercano.
- Renderiza en un solo `requestAnimationFrame` y solo cuando haya cambios (flag "dirty"); el cálculo del motor es memoizado con `useMemo`.

# 6. MOTOR MATEMÁTICO (`src/lib/physics.ts`, funciones puras, sin dependencias de React)
Para n bloques:
- Aᵢ = bᵢ·hᵢ; A = ΣAᵢ; M = Σmᵢ; ρᵢ = mᵢ/Aᵢ.
- Centroide geométrico: x̄_G = ΣAᵢxᵢ/ΣAᵢ, ȳ_G = ΣAᵢyᵢ/ΣAᵢ, con (xᵢ, yᵢ) los centros de los bloques.
- Centro de masa: x̄_CM = Σmᵢxᵢ/Σmᵢ, ȳ_CM = Σmᵢyᵢ/Σmᵢ.
- Estabilidad: xmin = min(xA,xB), xmax = max(xA,xB). Estable si xmin ≤ x̄_CM ≤ xmax. Margen = min(x̄_CM − xmin, xmax − x̄_CM) (negativo si cae fuera). Pivote crítico: A/B correspondiente a xmin si x̄_CM < xmin, o a xmax si x̄_CM > xmax.
- Reacciones (si es estable, con L = xmax − xmin > 0): R_min = M·g·(xmax − x̄_CM)/L, R_max = M·g·(x̄_CM − xmin)/L; verifica ΣF_y = 0 y Στ = 0 y muéstralo en el panel.
- Cada función exportada con tipos estrictos y documentada con comentario TSDoc. Devuelve también las "trazas" de cada paso (sumandos con valores) para que el panel pueda mostrar la sustitución numérica sin recalcular.
- El panel KaTeX muestra, en orden: (1) tabla de bloques con Aᵢ, mᵢ, ρᵢ, centros; (2) fórmula integral → sumatoria → sustitución numérica → resultado, para x̄_G, ȳ_G, x̄_CM, ȳ_CM; (3) criterio de estabilidad con los valores; (4) reacciones. Incluye una nota breve que explique por qué la integral de un bloque homogéneo se reduce a "área × centro".

# 7. SIMULACIÓN DE GRAVEDAD
- Si es ESTABLE: dibuja los vectores R_A y R_B (flechas verticales hacia arriba, longitud proporcional a la fuerza con escala visible y etiqueta en N) y el peso W = M·g aplicado en el CM.
- Si es INESTABLE: vuelco como cuerpo rígido alrededor del pivote crítico P:
  - Momento de inercia respecto al pivote (teorema de ejes paralelos): I_P = Σ [ mᵢ(bᵢ² + hᵢ²)/12 + mᵢ·dᵢ² ], con dᵢ la distancia del centro del bloque a P (usar metros y kg).
  - Dinámica: con r = CM − P rotado el ángulo θ, α = −M·g·(r_x cosθ − r_y sinθ)/I_P (signo de θ positivo antihorario). Integración semi-implícita de Euler (ω += α·dt; θ += ω·dt) con dt fijo = 1/120 s y acumulador de tiempo, para que sea independiente de los FPS.
  - Fin de la animación: cuando |θ| alcanza 90° o cuando un vértice de un bloque (excluyendo el pivote y los vértices que ya estaban en y = 0 en el lado del giro) toque el suelo y = 0. En el impacto, impacto inelástico: ω = 0 y estado "VOLCADO".
  - Dibuja todos los bloques, el CM y la plomada rotados con la misma transformación; la UI muestra θ, ω y t en vivo en el panel. "Reiniciar" restaura el estado inicial al instante. Mientras la simulación corre, la edición queda deshabilitada.

# 8. PRESETS (datos exactos, en `src/data/presets.ts`)
Cada preset: bloques `[b, h, m, x, y]`, apoyos y una tarjeta explicativa de 2–3 frases. Resultados esperados (se usan como tests).
1. **Torre Simétrica Estable**. A = −30, B = 30. Bloques: [60,10,20,−30,0], [30,30,15,−15,10], [20,20,8,−10,40]. Esperado: x̄_G = 0, x̄_CM = 0, M = 43, ESTABLE, R_A = R_B.
2. **Voladizo con Contrapeso**. A = −10, B = 30. Bloques: [40,10,20,−10,0], [50,10,10,20,10], [15,15,60,−10,10]. Esperado: x̄_G ≈ 23.06, x̄_CM ≈ 5.56, M = 90, ESTABLE.
3. **Colapso por Geometría**. A = −10, B = 10. Bloques de igual densidad (0.1 kg/cm²): [40,10,40,−10,0], [30,20,60,10,10]. Esperado: x̄_G = x̄_CM = 19, M = 100, INESTABLE, pivote B, margen = −9 cm.
4. **La Trampa de Densidad**. A = −20, B = 20. Bloques: [40,10,10,−20,0], [20,20,5,0,10], [20,10,80,20,0]. Esperado: x̄_G = 10 (dentro), x̄_CM ≈ 25.79 (fuera), M = 95, INESTABLE, pivote B. La tarjeta debe señalar explícitamente la diferencia entre centroide y centro de masa.

# 9. EXPORTACIÓN
- "Exportar PNG": captura del canvas a resolución ×2 con rótulo (nombre del preset, fecha, estado).
- "Reporte técnico": vista imprimible (`window.print()` con una hoja de estilos `@media print` en fondo blanco) con tabla de bloques, resultados, estado de estabilidad, el desarrollo KaTeX y la imagen del lienzo.

# 10. CALIDAD
- TypeScript estricto, sin `any`. Componentes pequeños y separados: `TopBar`, `BlockList`, `BlockInspector`, `SupportsInspector`, `Viewport`, `AnalysisPanel`, `StabilityCard`, `PresetCards`.
- Tests Vitest (`src/lib/physics.test.ts`): los 4 presets con tolerancia 1e-2, casos borde (M = 0, xA = xB, xA > xB), equilibrio de reacciones, e I_P de un bloque único contra la fórmula conocida.
- Rendimiento: sin re-renders del panel por frames del lienzo; el panel derecho solo se actualiza al cambiar el estado del modelo.
- README en español: requisitos (solo Docker), comandos, estructura del proyecto, explicación del modelo físico y despliegue en GitHub Pages.

# 11. CRITERIOS DE ACEPTACIÓN
- [ ] `docker compose up` sirve la app en el puerto 5173 y el HMR funciona.
- [ ] `docker compose run --rm app npm run build` termina sin errores y `dist/index.html` carga con rutas relativas (probado abriendo `dist` bajo una subruta).
- [ ] `docker compose run --rm app npm run test` pasa todos los tests.
- [ ] Los 4 presets reproducen los resultados esperados de la sección 8.
- [ ] El preset 4 muestra x̄_G dentro de la base y x̄_CM fuera, con vuelco animado en torno a B.
- [ ] Arrastrar bloques y apoyos actualiza panel, marcadores y plomada en tiempo real, sin tirones.
- [ ] No hay peticiones de red externas en producción (fuentes y KaTeX empaquetados).
- [ ] No hay NaN visibles ni errores en consola en los casos borde.

Comienza ahora únicamente con la **FASE 1**.
\end{Prompt}

\subsection{Correcciones e instrucciones posteriores}

Cada instrucción se reproduce literalmente, seguida del motivo y de lo que se obtuvo.

\subsubsection*{Corrección 1 — Sesión 2 (responsable: Hugo)}
\begin{Prompt}
Continua el desarrollo, aun esta lejos de ser un laboratorio de centroides y centros de masa
\end{Prompt}
\textbf{Motivo.} Tras la Fase 1 la aplicación sólo tenía la infraestructura y el esqueleto visual; todavía
no calculaba nada. \textbf{Resultado.} Fase 2: motor \texttt{physics.ts} con trazas de cada sumando,
presets como datos, panel con tarjeta de diagnóstico, reacciones y desarrollo KaTeX.

\subsubsection*{Corrección 2 — Sesión 3 (responsable: Daniela)}
\begin{Prompt}
continua
\end{Prompt}
\textbf{Resultado.} Fase 3: lienzo CAD interactivo. En la revisión se corrigieron dos defectos de fases
anteriores: el nombre de un bloque se recortaba con cada tecla (no se podían escribir espacios) y
«Centrar plano» forzaba el zoom al 100\,\% en lugar de encuadrar el modelo.

\subsubsection*{Corrección 3 — Sesión 4 (responsables: Sara y Hugo)}
\begin{Prompt}
continua
\end{Prompt}
\textbf{Resultado.} Fase 4: simulación de gravedad, tarjetas de presets, exportación PNG ($\times 2$ con
rótulo), reporte técnico imprimible y README.

\subsubsection*{Corrección 4 — Sesión 5 (responsable: Daniela)}
\begin{Prompt}
Continua hasta terminar
\end{Prompt}
\textbf{Resultado.} Entrega de los dos archivos restantes (\texttt{render.ts}, \texttt{Viewport.tsx}) y
cierre de los 8 criterios de aceptación.

\subsubsection*{Corrección 5 — Sesión 6 (responsables: los tres)}
\begin{Prompt}
Mejora la estetica visual, dale estilo enterprise, agrega para generar figuras mediante ecuaciones, drag and drop de figuras geometricas
\end{Prompt}
Antes de implementar, el agente planteó cuatro preguntas de diseño. Respuestas del equipo:

\begin{center}
\small
\begin{tabularx}{\textwidth}{L{5.2cm}Y}
\toprule
\textbf{Pregunta} & \textbf{Decisión del equipo}\\
\midrule
Tipo de figura por ecuación & Región entre $f(x)$ y $g(x)$ en $[a,b]$\\
Figuras de la paleta & Rectángulo y triángulo; círculo y semicírculo; cuarto de círculo y trapecio; polígono libre\\
Alcance del estilo \emph{enterprise} & Rediseño completo\\
Forma de entrega & Implementar y resumir\\
\bottomrule
\end{tabularx}
\end{center}

\textbf{Resultado.} Modelo generalizado a 8 tipos de figura, analizador de expresiones propio (sin
\texttt{eval}), integración numérica de Simpson, paleta con arrastrar y soltar, diálogos de ecuación y de
polígono, inspector por tipo de figura y nueva interfaz (barra de título, cinta de herramientas agrupada,
barra lateral de actividades, pestañas y barra de estado). Se añadieron los presets \emph{Arco Parabólico}
y \emph{Figuras Combinadas}.

\subsubsection*{Corrección 6 — Sesión 7 (responsables: Hugo, Daniela y Sara)}
\begin{Prompt}
Corrige un error con la simulacion, le doy a simular pero no hace nada por alguna razon, ademas generame el codigo Latex de una bitacora que contenga el desglose de esto realizado por 3 usuarios, Sara, Hugo y Daniela, desglosarlo a forma de que los 3 tuvieron trabajo uniforme, coloca el primer promt que te di y las correcciones que te fui dando para demostrar las evidencias en el desarrollo de la app
\end{Prompt}
\textbf{Resultado.} Ver la sección~\ref{sec:sim} y este documento.

% ---------------------------------------------------------------------
\section{Modelo matemático implementado}

\subsection{De la integral a la suma}
Para una región plana $R$ con densidad superficial $\rho$:
\[
\bar x_G=\frac{\iint_R x\,dA}{\iint_R dA},\qquad
\bar x_{CM}=\frac{\iint_R x\,\rho\,dA}{\iint_R \rho\,dA}.
\]
Para un rectángulo homogéneo $[x_0,x_0+b]\times[y_0,y_0+h]$:
\[
\iint x\,dA = h\,\frac{(x_0+b)^2-x_0^2}{2}= b\,h\left(x_0+\frac b2\right)=A_i\,x_i,
\]
y por la aditividad de la integral sobre regiones disjuntas,
$\bar x_G=\sum A_i x_i/\sum A_i$ y $\bar x_{CM}=\sum m_i x_i/\sum m_i$.

\subsection{Regiones definidas por ecuaciones}
Para la región $g(x)\le y\le f(x)$, $x\in[a,b]$ (regla de Simpson con $n=600$):
\[
A=\int_a^b\!\big[f(x)-g(x)\big]dx,\quad
\bar x=\frac1A\int_a^b\! x\big[f(x)-g(x)\big]dx,\quad
\bar y=\frac1A\int_a^b\!\tfrac12\big[f(x)^2-g(x)^2\big]dx .
\]
Comprobación (preset \emph{Arco Parabólico}): $f(x)=20-\tfrac{x^2}{20}$, $g=0$, $[-20,20]$
$\Rightarrow A=533{,}33\ \text{cm}^2$, $\bar y=8\ \text{cm}=\tfrac25 h$.

\subsection{Estabilidad, reacciones y vuelco}
\begin{itemize}
  \item Estable si $x_{\min}\le \bar x_{CM}\le x_{\max}$; margen $=\min(\bar x_{CM}-x_{\min},\,x_{\max}-\bar x_{CM})$.
  \item $R_{\min}=W\,\dfrac{x_{\max}-\bar x_{CM}}{L}$, \quad $R_{\max}=W\,\dfrac{\bar x_{CM}-x_{\min}}{L}$, \quad $W=Mg$.
  \item Vuelco: $I_P=\sum\big[I_{c,i}+m_i d_i^2\big]$, \quad
        $\alpha=-\dfrac{Mg\,(r_x\cos\theta-r_y\sin\theta)}{I_P}$, Euler semi-implícito con $\Delta t=1/120$ s.
\end{itemize}

% ---------------------------------------------------------------------
\section{Verificación y pruebas}

\subsection{Evolución de las pruebas automatizadas}
\begin{center}
\small
\begin{tabular}{lrl}
\toprule
\textbf{Etapa} & \textbf{Tests} & \textbf{Contenido principal}\\
\midrule
Fase 1 & 8 & Reducer del modelo y validación de rangos\\
Fase 2 & 30 & Presets 1–4, casos borde, equilibrio, $I_P$\\
Fase 3 & 38 & Cámara, cuadrícula adaptativa, sombreado por densidad\\
Fase 4 & 46 & Dinámica del vuelco, independencia de FPS, energía\\
Figuras y ecuaciones & 88 & Fórmulas exactas por figura, analizador de expresiones, Simpson\\
Interfaz \emph{enterprise} & 103 & Fábrica de figuras (posición y validez de cada tipo)\\
\bottomrule
\end{tabular}
\end{center}

\subsection{Resultados de los presets}
\begin{center}
\small
\begin{tabular}{llllll}
\toprule
\textbf{Preset} & $\bar x_G$ & $\bar x_{CM}$ & $M$ & \textbf{Estado} & \textbf{Comprobación}\\
\midrule
1. Torre Simétrica & 0 & 0 & 43 kg & Estable & $R_A=R_B=210{,}92$ N \ok\\
2. Voladizo con Contrapeso & 23,06 & 5,56 & 90 kg & Estable & \ok\\
3. Colapso por Geometría & 19 & 19 & 100 kg & Inestable (B) & margen $-9$ cm \ok\\
4. Trampa de Densidad & 10 & 25,79 & 95 kg & Inestable (B) & vuelco a $-26{,}57^\circ$ \ok\\
5. Arco Parabólico & 0 & 0 & 30 kg & Estable & $\bar y_G = 8$ cm \ok\\
6. Figuras Combinadas & 3,29 & 3,89 & 60 kg & Estable & \ok\\
\bottomrule
\end{tabular}
\end{center}

\subsection{Criterios de aceptación}
\begin{center}
\small
\begin{tabularx}{\textwidth}{Yc}
\toprule
\textbf{Criterio} & \textbf{Estado}\\
\midrule
\texttt{docker compose up} sirve la app en el puerto 5173 y el HMR funciona & \ok\\
El build termina sin errores y carga con rutas relativas bajo una subruta & \ok\\
Todos los tests pasan & \ok\\
Los presets reproducen los resultados esperados & \ok\\
Preset 4: $\bar x_G$ dentro, $\bar x_{CM}$ fuera y vuelco animado en torno a B & \ok\\
El arrastre actualiza panel, marcadores y plomada en tiempo real & \ok\\
Sin peticiones de red externas (fuentes y KaTeX empaquetados) & \ok\\
Sin NaN visibles ni errores en consola en los casos borde & \ok\\
\bottomrule
\end{tabularx}
\end{center}

\subsection{Defectos detectados y corregidos}
\begin{center}
\small
\begin{tabularx}{\textwidth}{L{2.4cm}YY}
\toprule
\textbf{Etapa} & \textbf{Defecto} & \textbf{Corrección}\\
\midrule
Fase 3 & El nombre del bloque se recortaba en cada tecla & El recorte sólo se aplica al validar un nombre vacío\\
Fase 3 & «Centrar plano» forzaba el zoom al 100\,\% & Ahora encuadra el modelo completo\\
Figuras & Error de signo en el recorte de polígonos (área de solapamiento) & Parámetro de intersección corregido\\
Figuras & Conversión de unidades aplicada dos veces en $I_P$ & $I_P=\sum m_i J_{c,i}/(A_i\cdot 10^4)+m_i d_i^2$\\
Interfaz & La barra de estado no mostraba VOLCADO & Suscripción al estado de la simulación\\
Interfaz & Etiquetas de la cinta recortadas y títulos duplicados & Botones más anchos y títulos únicos\\
Simulación & «Le doy a simular pero no hace nada» & Ver sección~\ref{sec:sim}\\
\bottomrule
\end{tabularx}
\end{center}

% ---------------------------------------------------------------------
\section{Corrección del error de simulación}\label{sec:sim}

\subsection{Diagnóstico}
Se reprodujo el problema en el navegador con distintos modelos y se midió cuántos fotogramas distintos
dibujaba el lienzo durante 3 segundos. La simulación sí se ejecutaba, pero en varios casos no se percibía:

\begin{enumerate}
  \item \textbf{Vuelco demasiado rápido.} En los presets 3 y 4 el vuelco completo dura unos 0,19 s
        simulados, que en tiempo real equivalen a unos pocos fotogramas.
  \item \textbf{Giro imperceptible en figuras curvas.} Los contornos de círculos, semicírculos y regiones
        por ecuación se muestrean con muchos vértices; los vértices vecinos al punto de apoyo quedan a
        fracciones de centímetro del suelo. Si la figura descansa sobre el suelo más allá del pivote, el
        primer impacto llega tras un giro de apenas $0{,}06^\circ$: físicamente el suelo impide el vuelco,
        pero la pantalla no cambia y la tarjeta pasa directamente a VOLCADO.
  \item \textbf{Modo estable estático.} En modelos estables sólo se añadían flechas de fuerza, sin
        movimiento ni mensaje.
  \item \textbf{Atajo de teclado inexistente.} La tecla \texttt{G} aparecía en la ayuda pero no estaba conectada.
\end{enumerate}

\subsection{Solución}
La regla física de contacto no se modificó (es correcta). Se hizo la simulación perceptible y explicativa:
\begin{itemize}
  \item cámara lenta configurable ($\times1$, $\times0{,}5$, $\times0{,}25$, $\times0{,}1$; por defecto
        $\times0{,}25$) sin alterar el paso fijo de integración, y botón \emph{Repetir};
  \item aviso explicativo cuando el giro máximo es menor de $3^\circ$, indicando el vértice que toca el suelo y
        cómo obtener un vuelco mayor;
  \item marca de \emph{impacto} en el vértice que toca el suelo;
  \item en modo estable, mensaje «sistema estable: no hay movimiento» y animación de crecimiento de las
        flechas $R_A$, $R_B$ y $W$;
  \item atajo \texttt{G} para iniciar y detener la simulación.
\end{itemize}

% ---------------------------------------------------------------------
\section{Estructura final del proyecto}
\begin{Prompt}
├── Dockerfile, docker-compose.yml     Entorno de desarrollo y build sin dependencias locales
├── .github/workflows/deploy.yml       Despliegue a GitHub Pages
├── bitacora/bitacora.tex              Esta bitácora
└── src/
    ├── lib/        physics.ts, shapes.ts, expr.ts, dynamics.ts, derivation.ts, geometry.ts,
    │               camera.ts, density.ts, shapeFactory.ts (+ tests *.test.ts)
    ├── data/       presets.ts (6 presets con resultados esperados)
    ├── canvas/     render.ts (dibujo 2D nativo)
    ├── state/      modelo, análisis, vista, simulación, herramientas, reporte
    └── components/ TitleBar, Ribbon, LeftPanel, RightPanel, StatusBar, Viewport, Palette,
                    ShapeInspector, EquationDialog, PolygonDialog, SimulationCard, ...
\end{Prompt}

% ---------------------------------------------------------------------
\section{Conclusiones}

\begin{itemize}
  \item Especificar con precisión (unidades, convenciones, datos de prueba y criterios de aceptación)
        permitió verificar cada fase de forma objetiva: los resultados esperados de los presets se
        convirtieron en pruebas automatizadas.
  \item El trabajo por fases, con revisión antes de cada «continúa», ayudó a detectar defectos pronto.
  \item Las correcciones más valiosas surgieron de usar la aplicación como lo haría un estudiante:
        el error de la simulación no era un fallo del cálculo sino de cómo se presentaba.
  \item La aplicación conecta directamente el contenido del curso: la integral
        $\iint_R x\,\rho\,dA$ se reduce a sumas para figuras simples y se calcula numéricamente para
        regiones entre curvas.
\end{itemize}

\end{document}
